Nullniveau der Lageenergie: der Abschussbolzen. Zustand 1: die Feder gespannt (nur Spannenergie). Zustand 2: die Kugel verlässt die Feder (nur kinetische Energie). Zustand 3: zuoberst auf dem Tisch (Lageenergie und kinetische Energie). Zuerst die Einheiten: Zentimeter in Meter, Gramm in Kilogramm.

\begin{abcliste}
\abc
Sei $k$ die Federkonstante, $s$, wie weit sie zurückgezogen wird, und $m$ die Masse der Kugel. Die Spannenergie im Zustand 1 wird zur kinetischen Energie im Zustand 2:
\begin{align}
 \frac{1}{2}\,k\,s^2 &= \frac{1}{2}\,m\,v_1^2 \\
 v_1 &= \vaF \\
   &= \sF\cdot\sqrt{\frac{\k}{\m}} \\
   &= \va \approx \result{\vaP{0}{2}}
\end{align}

\abc
Sei $h$ die Höhe des oberen Endes über dem Bolzen. Ein Teil der Spannenergie ist zu Lageenergie geworden; der Rest ist noch kinetische Energie. Zwischen den Zuständen 1 und 3:
\begin{align}
 \frac{1}{2}\,k\,s^2 &= m\,g\,h + \frac{1}{2}\,m\,v_2^2 \\
 v_2 &= \vbF \\
   &= \sqrt{\frac{\k\cdot\left(\sF\right)^2}{\m}-2\cdot\g\cdot\hT} \\
   &= \vb \approx \result{\vbP{0}{2}}
\end{align}
Geschwindigkeiten kann man nicht subtrahieren, Energien schon.
\end{abcliste}
